\part{Universal inversion calculus}

This part collects the inversion machinery used throughout the volume.  Its
purpose is not merely to quote familiar formulas: coefficient conventions,
domains, branches, and certification hypotheses are made explicit so that
later applications can invoke the results without changing their meaning.
Unless otherwise stated, a power series is based at the origin and
\([z^n]A(z)\) denotes its \(z^n\)-coefficient.

\section{Regular inverse germs}
\label{sec:uq-regular-germs}

\begin{definition}[Partial inverse]
\label{def:uq-partial-inverse}
Let \(F(p,\lambda)\) be a scalar observable, where \(p\) is the variable to be
recovered and \(\lambda=(\lambda_1,\ldots,\lambda_s)\) consists of passive
parameters.  A \emph{partial inverse germ in \(p\)} at
\((p_0,\lambda_0)\), with \(y_0=F(p_0,\lambda_0)\), is a germ
\(G(y,\lambda)\) at \((y_0,\lambda_0)\) satisfying
\[
 F(G(y,\lambda),\lambda)=y,
 \qquad
 G(y_0,\lambda_0)=p_0.
\]
Thus the word \emph{inverse} always includes the recovered variable, the
base point, and the selected local branch.
\end{definition}

\begin{theorem}[Regular analytic partial inverse]
\label{thm:uq-regular-analytic-inverse}
Let \(F\) be holomorphic in a neighbourhood of
\((p_0,\lambda_0)\in\ComplexNumbers\times\ComplexNumbers^s\), put
\(y_0=F(p_0,\lambda_0)\), and suppose
\[
 F_p(p_0,\lambda_0)\ne0.
\]
Then there are neighbourhoods \(U\) of \((y_0,\lambda_0)\) and \(V\) of
\(p_0\), and a unique holomorphic map \(G\colon U\to V\), after shrinking
them if necessary, which is the partial inverse germ of
\cref{def:uq-partial-inverse}.  For each fixed nearby \(\lambda\), the two
maps \(p\mapsto F(p,\lambda)\) and \(y\mapsto G(y,\lambda)\) are inverse
biholomorphic germs.
\end{theorem}

\begin{proof}
Apply the holomorphic implicit-function theorem to
\[
 \Psi(p,y,\lambda):=F(p,\lambda)-y.
\]
At \((p_0,y_0,\lambda_0)\) one has \(\Psi=0\) and
\(\Psi_p=F_p\ne0\), so the theorem supplies a unique holomorphic solution
\(p=G(y,\lambda)\).  Its defining identity is
\(F(G(y,\lambda),\lambda)=y\).  Holding \(\lambda\) fixed and applying the
same theorem with \(p\) as input shows that this identity holds between
inverse germs in both directions.  Uniqueness is the uniqueness clause of
the implicit-function theorem.
\end{proof}

\begin{remark}
The identical statement holds in the real \(C^r\) category, \(1\le r\le
\infty\), with \emph{holomorphic} replaced by \(C^r\).  This is the
one-dimensional parameterized inverse-function theorem.
\end{remark}

\section{Formal triangular inversion}
\label{sec:uq-formal-inversion}

\begin{theorem}[Triangular inverse recurrence]
\label{thm:uq-triangular-inverse}
Let \(R\) be a commutative ring and
\[
 f(z)=\sum_{n\ge1}c_nz^n\in R[[z]]
\]
with \(c_1\) a unit.  There is a unique series
\(g(w)=\sum_{n\ge1}b_nw^n\) satisfying both
\[
 f(g(w))=w,\qquad g(f(z))=z.
\]
Its coefficients are determined without division except by \(c_1\):
\[
 b_1=c_1^{-1},
\]
and, for \(n\ge2\),
\begin{equation}
\label{eq:uq-triangular-recurrence}
 b_n=-c_1^{-1}\sum_{k=2}^{n}c_k
 \sum_{\substack{j_1+\cdots+j_k=n\\j_1,\ldots,j_k\ge1}}
 b_{j_1}\cdots b_{j_k}.
\end{equation}
\end{theorem}

\begin{proof}
The coefficient of \(w\) in \(f(g(w))\) is \(c_1b_1\), which forces
\(b_1=c_1^{-1}\).  For \(n\ge2\), the coefficient of \(w^n\) is
\[
 c_1b_n+
 \sum_{k=2}^{n}c_k
 \sum_{\substack{j_1+\cdots+j_k=n\\j_i\ge1}}
 b_{j_1}\cdots b_{j_k}.
\]
Every \(j_i\) in the second term is at most \(n-1\), so setting this
coefficient equal to zero gives \eqref{eq:uq-triangular-recurrence} and
constructs a unique right compositional inverse \(g\).

It remains to check the other composition.  Put \(h(z)=g(f(z))\).  Then
\[
 f(h(z))=(f\circ g)(f(z))=f(z),
\]
and the linear coefficient of \(h\) is \(b_1c_1=1\).  If
\(h(z)\ne z\), let \(r\ge2\) be the first degree at which they differ and
write \(h(z)=z+dz^r+\BigO(z^{r+1})\) with \(d\ne0\).  The first degree at
which \(f(h(z))\) and \(f(z)\) differ is then \(r\), with difference
\(c_1d\).  Since \(c_1\) is a unit this is nonzero, contradicting
\(f\circ h=f\).  Hence \(g\circ f\) is also the identity.
\end{proof}

\begin{corollary}[The first five ordinary coefficients]
\label{cor:uq-first-inverse-coefficients}
Over a field, the first coefficients in
\cref{thm:uq-triangular-inverse} are
\[
\begin{aligned}
 b_1&=\frac1{c_1},&
 b_2&=-\frac{c_2}{c_1^3},\\
 b_3&=\frac{2c_2^2-c_1c_3}{c_1^5},&
 b_4&=\frac{-5c_2^3+5c_1c_2c_3-c_1^2c_4}{c_1^7},\\
 b_5&=\frac{14c_2^4-21c_1c_2^2c_3
       +6c_1^2c_2c_4+3c_1^2c_3^2-c_1^3c_5}{c_1^9}.
\end{aligned}
\]
\end{corollary}

\begin{proof}
The recurrence specializes successively to
\[
\begin{aligned}
b_2={}&-c_1^{-1}c_2b_1^2,\\
b_3={}&-c_1^{-1}(2c_2b_1b_2+c_3b_1^3),\\
b_4={}&-c_1^{-1}\{c_2(2b_1b_3+b_2^2)
                    +3c_3b_1^2b_2+c_4b_1^4\},\\
b_5={}&-c_1^{-1}\{c_2(2b_1b_4+2b_2b_3)
 +c_3(3b_1^2b_3+3b_1b_2^2)+4c_4b_1^3b_2+c_5b_1^5\}.
\end{aligned}
\]
Substituting \(b_1=c_1^{-1}\) from top to bottom and collecting terms gives
the five displayed formulas.
\end{proof}

\begin{remark}
These are ordinary-power-series coefficients.  Factorial-normalized
derivatives obey the Bell recurrence in
\cref{thm:uq-bell-inverse-recursion}.
\end{remark}

\section{Formal and analytic Lagrange--B\"urmann inversion}
\label{sec:uq-lagrange-burmann}

\begin{lemma}[Formal residue under a change of coordinate]
\label{lem:uq-formal-residue-substitution}
Let \(R\) be a commutative \(\RationalNumbers\)-algebra, let
\(f(z)=zU(z)\) with \(U(0)\) a unit, and let
\(\operatorname{res}_z A(z):=[z^{-1}]A(z)\).  For every formal Laurent
series \(A(w)\in R((w))\),
\[
 \operatorname{res}_w A(w)
 =
 \operatorname{res}_z\!\bigl(A(f(z))f'(z)\bigr).
\]
\end{lemma}

\begin{proof}
Both sides are \(R\)-linear, and in any requested coefficient only finitely
many terms of \(A\) contribute, so it suffices to take \(A(w)=w^k\).
If \(k\ne-1\), then
\[
 f(z)^kf'(z)=\frac1{k+1}\frac{\Differential}{\Differential z}f(z)^{k+1},
\]
whose residue is zero, as is the residue of \(w^k\).  Division by \(k+1\)
is legitimate because \(R\) is a \(\RationalNumbers\)-algebra.  If \(k=-1\), then
\[
 \frac{f'(z)}{f(z)}=\frac1z+\frac{U'(z)}{U(z)}.
\]
The second summand is an ordinary power series and has zero residue, while
the first has residue one.  This is also
\(\operatorname{res}_w w^{-1}\).
\end{proof}

\begin{theorem}[Formal Lagrange--B\"urmann formula]
\label{thm:uq-formal-lagrange-burmann}
Let \(R\) be a commutative \(\RationalNumbers\)-algebra, let
\(\phi(z)\in R[[z]]\) have invertible constant term, and put
\[
 f(z)=\frac{z}{\phi(z)}.
\]
Let \(g\) be the compositional inverse of \(f\).  Then for every
\(H(z)\in R[[z]]\) and every \(n\ge1\),
\begin{equation}
\label{eq:uq-lagrange-burmann}
 [w^n]H(g(w))
 =\frac1n[z^{n-1}]H'(z)\phi(z)^n.
\end{equation}
In particular,
\[
 [w^n]g(w)=\frac1n[z^{n-1}]\phi(z)^n.
\]
\end{theorem}

\begin{proof}
By \cref{lem:uq-formal-residue-substitution}, with \(w=f(z)\),
\[
\begin{aligned}
 [w^n]H(g(w))
 &=\operatorname{res}_w H(g(w))w^{-n-1}\\
 &=\operatorname{res}_z H(z)f'(z)f(z)^{-n-1}.
\end{aligned}
\]
The residue of a formal derivative is zero.  Applying this to
\(H(z)f(z)^{-n}\) gives
\[
 0=\operatorname{res}_z\!\left(
 H'(z)f(z)^{-n}-nH(z)f'(z)f(z)^{-n-1}\right).
\]
Consequently the preceding residue equals
\[
 \frac1n\operatorname{res}_z H'(z)f(z)^{-n}
 =\frac1n\operatorname{res}_z
 H'(z)z^{-n}\phi(z)^n,
\]
which is \eqref{eq:uq-lagrange-burmann}.  Taking \(H(z)=z\) proves the
special case.
\end{proof}

\begin{theorem}[Analytic Lagrange--B\"urmann formula]
\label{thm:uq-analytic-lagrange}
Under the hypotheses of \cref{thm:uq-regular-analytic-inverse} with the
passive parameter fixed, write \(y_0=F(p_0)\).  The selected inverse germ is
\[
 G(y)=p_0+\sum_{n\ge1}\frac{(y-y_0)^n}{n!}
 \left[
 \frac{\Differential^{\,n-1}}{\Differential p^{\,n-1}}
 \left(\frac{p-p_0}{F(p)-F(p_0)}\right)^{\!n}
 \right]_{p=p_0}.
\]
The quotient in parentheses is understood by its removable value
\(1/F'(p_0)\).
\end{theorem}

\begin{proof}
Set \(z=p-p_0\), \(w=y-y_0\), and
\[
 f(z)=F(p_0+z)-F(p_0),\qquad \phi(z)=\frac{z}{f(z)}.
\]
Since \(f'(0)=F'(p_0)\ne0\), \(\phi\) is holomorphic at zero with nonzero
constant term.  The analytic inverse-function theorem gives a convergent
inverse germ \(z=g(w)\), while
\cref{thm:uq-formal-lagrange-burmann} gives
\[
 [w^n]g(w)=\frac1n[z^{n-1}]\phi(z)^n
 =\frac1{n!}\left[\frac{\Differential^{\,n-1}}{\Differential z^{\,n-1}}
 \phi(z)^n\right]_{z=0}.
\]
Translating back from \(z\) and \(w\) proves the formula.
\end{proof}

\begin{warning}[Radius is branchwise]
\label{warn:uq-inverse-radius}
Lagrange--B\"urmann determines a germ, not its global sheet.  Its radius of
convergence is the distance from the expansion value to the nearest
singularity of the \emph{selected analytically continued inverse germ}.
A critical value on another sheet, or a value reached only after crossing a
cut, need not control that radius.  Thus the phrase \emph{the nearest
critical value} is not a valid unqualified radius formula.
\end{warning}

\section{Bell-polynomial packaging}
\label{sec:uq-bell-packaging}

The complete and partial exponential Bell polynomials are fixed by
\begin{align}
 \exp\!\left(\sum_{j\ge1}x_j\frac{t^j}{j!}\right)
 &=\sum_{n\ge0}B_n(x_1,\ldots,x_n)\frac{t^n}{n!},
 \label{eq:uq-complete-bell-definition}\\
 \frac1{r!}\left(\sum_{j\ge1}x_j\frac{t^j}{j!}\right)^{\!r}
 &=\sum_{n\ge r}B_{n,r}(x_1,\ldots,x_{n-r+1})
   \frac{t^n}{n!}.
 \label{eq:uq-partial-bell-definition}
\end{align}
These identities are definitions and fix all factorial conventions.

\begin{lemma}[Exponentiating logarithmic jets]
\label{lem:uq-bell-exponentiation}
Let \(R\) be a commutative \(\RationalNumbers\)-algebra and let \(A(z)\in R[[z]]\) have
invertible constant term.  If
\[
 \log\frac{A(z)}{A(0)}
 =\sum_{j\ge1}\ell_j\frac{z^j}{j!},
\]
then
\[
 \frac{A(z)}{A(0)}
 =\sum_{n\ge0}B_n(\ell_1,\ldots,\ell_n)\frac{z^n}{n!}.
\]
\end{lemma}

\begin{proof}
Exponentiate the displayed logarithmic identity.  The right-hand side is
exactly the defining generating function
\eqref{eq:uq-complete-bell-definition}.
\end{proof}

\begin{theorem}[Bell recurrence for inverse derivatives]
\label{thm:uq-bell-inverse-recursion}
Let \(R\) be a commutative \(\RationalNumbers\)-algebra and write
\[
 f(z)=\sum_{r\ge1}f_r\frac{z^r}{r!},\qquad
 g(w)=\sum_{n\ge1}g_n\frac{w^n}{n!},
\]
where \(f_1\) is a unit and \(f(g(w))=w\).  Then
\[
 g_1=f_1^{-1}
\]
and, for \(n\ge2\),
\begin{equation}
\label{eq:uq-bell-inverse-recurrence}
 g_n=-f_1^{-1}\sum_{r=2}^{n}
 f_rB_{n,r}(g_1,\ldots,g_{n-r+1}).
\end{equation}
\end{theorem}

\begin{proof}
Substituting \(g(w)\) into the \(r\)-th term of \(f\), and using
\eqref{eq:uq-partial-bell-definition}, shows that the coefficient multiplying
\(w^n/n!\) in \(f(g(w))\) is
\[
 \sum_{r=1}^{n}f_r
 B_{n,r}(g_1,\ldots,g_{n-r+1}).
\]
Here \(B_{n,1}(g_1,\ldots,g_n)=g_n\).  For \(n=1\), comparison with
\(w\) gives \(f_1g_1=1\).  For every \(n\ge2\), the displayed coefficient
must vanish; isolating \(f_1g_n\) gives
\eqref{eq:uq-bell-inverse-recurrence}.  Each right-hand side involves only
\(g_j\) with \(j<n\), so the recurrence is triangular.
\end{proof}

\section{Passive-parameter sensitivities}
\label{sec:uq-sensitivities}

\begin{theorem}[First and second inverse sensitivities]
\label{thm:uq-inverse-sensitivities}
Let \(F\) and its partial inverse \(G\) be \(C^2\) (or holomorphic) and suppose
\(F_p\ne0\) on the germ under consideration.  All derivatives of \(F\) below
are evaluated at \((G(y,\lambda),\lambda)\).  Then
\[
 G_y=\frac1{F_p},
 \qquad
 G_{\lambda_i}=-\frac{F_{\lambda_i}}{F_p},
\]
\[
 G_{yy}=-\frac{F_{pp}}{F_p^3},
 \qquad
 G_{y\lambda_i}
 =\frac{F_{pp}F_{\lambda_i}-F_pF_{p\lambda_i}}{F_p^3},
\]
and
\[
 G_{\lambda_i\lambda_j}
 =-\frac{
 F_{pp}G_{\lambda_i}G_{\lambda_j}
 +F_{p\lambda_i}G_{\lambda_j}
 +F_{p\lambda_j}G_{\lambda_i}
+F_{\lambda_i\lambda_j}}{F_p}.
\]
If \(F\) and \(G\) are \(C^r\), the same differentiation gives all mixed
derivatives through total order \(r\) recursively: at every new order the
highest derivative of \(G\) occurs once, multiplied by \(F_p\), and all
remaining terms involve lower derivatives of \(G\).  In the holomorphic case
this continues to every order.
\end{theorem}

\begin{proof}
Differentiate \(F(G(y,\lambda),\lambda)=y\) with respect to \(y\) and to
\(\lambda_i\):
\[
 F_pG_y=1,\qquad F_pG_{\lambda_i}+F_{\lambda_i}=0.
\]
These give the first two formulas.  Differentiating the first identity once
more with respect to \(y\) gives
\[
 F_{pp}G_y^2+F_pG_{yy}=0,
\]
and substituting \(G_y=1/F_p\) gives the formula for \(G_{yy}\).
Differentiating \(F_pG_y=1\) with respect to \(\lambda_i\) at fixed \(y\)
gives
\[
 (F_{pp}G_{\lambda_i}+F_{p\lambda_i})G_y
 +F_pG_{y\lambda_i}=0.
\]
Substitution of the first derivatives yields the stated mixed formula.
Finally, differentiate
\(F_pG_{\lambda_i}+F_{\lambda_i}=0\) with respect to
\(\lambda_j\).  The chain rule gives
\[
 F_pG_{\lambda_i\lambda_j}
 +F_{pp}G_{\lambda_i}G_{\lambda_j}
 +F_{p\lambda_j}G_{\lambda_i}
 +F_{p\lambda_i}G_{\lambda_j}
 +F_{\lambda_i\lambda_j}=0,
\]
which proves the last formula.  Repeating the chain rule proves the
triangular assertion at every order.
\end{proof}

\begin{remark}[Conditioning]
\label{rem:uq-inverse-conditioning}
For perturbations of the observed value with passive parameters fixed, the
absolute local condition number is \(\lvert F_p\rvert^{-1}\).  When
\(p\ne0\) and \(y\ne0\), the corresponding relative condition number is
\[
 \left\lvert\frac{y}{pF_p}\right\rvert.
\]
The blow-up as \(F_p\to0\) is a conditioning statement; by itself it is not
a proof of nonidentifiability.
\end{remark}

\section{Finite ramification and Puiseux inversion}
\label{sec:uq-puiseux}

\begin{theorem}[Finite-ramification normal form]
\label{thm:uq-finite-ramification}
Let \(F\) be holomorphic near \(p_0\), let \(y_0=F(p_0)\), and suppose, for
some \(m\ge2\),
\[
 F(p_0+z)
 =y_0+c_mz^m+c_{m+1}z^{m+1}+\BigO(z^{m+2}),
 \qquad c_m\ne0.
\]
Choose a local value
\[
 \xi=\left(\frac{y-y_0}{c_m}\right)^{1/m}
\]
on a sufficiently small simply connected punctured \(y\)-domain.  Then the
\(m\) inverse branches near \(p_0\), indexed by the \(m\)-th roots of unity,
are
\begin{equation}
\label{eq:uq-finite-ramification}
 p-p_0
 =\omega\xi-\frac{c_{m+1}}{mc_m}\omega^2\xi^2
 +\BigO(\xi^3),
 \qquad \omega^m=1.
\end{equation}
For every sufficiently small \(y\ne y_0\), these are precisely the \(m\)
solutions near \(p_0\), counted without multiplicity; at \(y=y_0\) they
coalesce.
\end{theorem}

\begin{proof}
Factor
\[
 F(p_0+z)-y_0=c_mz^mH(z),
 \qquad
 H(z)=1+\frac{c_{m+1}}{c_m}z+\BigO(z^2).
\]
On a small disk, \(H\) is nonzero and has a holomorphic \(m\)-th root with
value one at zero.  Define \(v=zH(z)^{1/m}\).  Then \(v(0)=0\) and
\(v'(0)=1\), so \(v\) has a regular holomorphic inverse \(z=\psi(v)\).
Moreover,
\[
 \frac{F(p_0+z)-y_0}{c_m}=v^m.
\]
Thus the solutions are \(z=\psi(\omega\xi)\), \(\omega^m=1\).  Since
\[
 v=z+\frac{c_{m+1}}{mc_m}z^2+\BigO(z^3),
 \qquad
 \psi(v)=v-\frac{c_{m+1}}{mc_m}v^2+\BigO(v^3),
\]
substitution gives \eqref{eq:uq-finite-ramification}.  For nonzero
\(\xi\), the \(m\) values \(\omega\xi\) are distinct, and the injectivity of
\(\psi\) on a smaller disk makes their images distinct.  Conversely every
nearby solution yields one of these \(m\) roots of \(v^m=\xi^m\), proving
exhaustion.
\end{proof}

\begin{corollary}[Universal fold jet through order two]
\label{cor:uq-fold-jet}
Let a holomorphic germ, in the notation of
\cref{thm:uq-finite-ramification} with \(m=2\), have the expansion
\[
 y-y_*=Az^2+Bz^3+Cz^4+Dz^5+\BigO(z^6),
 \qquad A\ne0,
\]
and choose a square root \(\sqrt A\).  The two inverse branches, indexed by
\(\sigma\in\{1,-1\}\), are
\begin{align}
 z={}&\sigma\frac{(y-y_*)^{1/2}}{\sqrt A}
 -\frac{B}{2A^2}(y-y_*)\notag\\
 &+\sigma\frac{5B^2-4AC}{8A^{7/2}}(y-y_*)^{3/2}
 -\frac{A^2D-3ABC+2B^3}{2A^5}(y-y_*)^2
 +\BigO((y-y_*)^{5/2}).
\label{eq:uq-fold-jet}
\end{align}
The powers \(A^{7/2}\) use the same chosen square root as the leading term.
\end{corollary}

\begin{proof}
Put \(t=(y-y_*)^{1/2}\) on the chosen ramified coordinate and seek
\[
 z=at+bt^2+ct^3+dt^4+\BigO(t^5).
\]
Matching the coefficient of \(t^2\) gives \(Aa^2=1\), hence
\(a=\sigma/\sqrt A\).  Successively setting the coefficients of
\(t^3,t^4,t^5\) equal to zero gives
\[
\begin{aligned}
 2Aab+Ba^3&=0,\\
 A(b^2+2ac)+3Ba^2b+Ca^4&=0,\\
 A(2ad+2bc)+B(3a^2c+3ab^2)+4Ca^3b+Da^5&=0.
\end{aligned}
\]
Solving these triangular equations yields
\[
 b=-\frac{B}{2A^2},\qquad
 c=\sigma\frac{5B^2-4AC}{8A^{7/2}},\qquad
 d=-\frac{A^2D-3ABC+2B^3}{2A^5}.
\]
Substitution into the ansatz proves \eqref{eq:uq-fold-jet}.
\end{proof}

\begin{warning}[What a singular Jacobian does not say]
\label{warn:uq-singular-jacobian}
The condition \(F_p(p_0,\lambda_0)=0\) says only that the regular
inverse-function theorem is unavailable.  A finite first nonzero derivative
can produce finitely ramified branches as above; an injective real map such
as \(p\mapsto p^3\) remains identifiable despite its zero derivative; and
more degenerate behaviour is also possible.  Neither nonidentifiability nor
Puiseux behaviour follows from a singular Jacobian alone.
\end{warning}

\section{Asymptotic inverse transfer and Lambert \(W\)}
\label{sec:uq-asymptotic-transfer}

\begin{proposition}[Residual transfer along a monotone inverse branch]
\label{prop:uq-asymptotic-residual-transfer}
Let \(F\) be real \(C^1\) and injective on an interval \(I\), with inverse
\(G\) on \(F(I)\).  Let \(\widetilde G(y)\in I\).  If
\(\lvert F'(p)\rvert\ge m(y)>0\) on the segment joining
\(\widetilde G(y)\) to \(G(y)\), then
\[
 \lvert\widetilde G(y)-G(y)\rvert
 \le
 \frac{\lvert F(\widetilde G(y))-y\rvert}{m(y)}.
\]
In particular, a forward residual estimate transfers to an inverse
asymptotic estimate once the derivative scale is controlled on the selected
branch.
\end{proposition}

\begin{proof}
The mean-value theorem supplies a point \(\theta\) between
\(\widetilde G(y)\) and \(G(y)\) such that
\[
 F(\widetilde G(y))-F(G(y))
 =F'(\theta)(\widetilde G(y)-G(y)).
\]
Use \(F(G(y))=y\) and the assumed lower bound for
\(\lvert F'(\theta)\rvert\).
\end{proof}

\begin{theorem}[Reciprocal-scale reversion]
\label{thm:uq-reciprocal-log-reversion}
Let \(A\ne0\) and suppose, on a real or complex sector as \(t\to0\),
\[
 x=X(t)=\frac{A}{t}+B+Ct+Dt^2+\BigO(t^3).
\]
For any inverse branch \(t=T(x)\to0\), put \(S=x-B\).  Then
\(t\sim A/S\) and, in the corresponding sector as \(S\to\infty\),
\begin{equation}
\label{eq:uq-reciprocal-reversion}
 T(x)=\frac{A}{S}+\frac{CA^2}{S^3}
 +\frac{DA^3}{S^4}+\BigO(S^{-5}).
\end{equation}
In particular, the coefficient of \(S^{-2}\) is zero.
\end{theorem}

\begin{proof}
After subtracting \(B\) and multiplying by \(t\), the inverse equation is
\[
 St=A+Ct^2+Dt^3+\BigO(t^4).
\]
Since \(t\to0\), this first gives \(St\to A\), hence \(t\sim A/S\).
Dividing the same identity by \(S\) and using \(t=\BigO(S^{-1})\) gives
\[
 t=\frac AS+\BigO(S^{-3}).
\]
It follows that
\[
 t^2=\frac{A^2}{S^2}+\BigO(S^{-4}),
 \qquad
 t^3=\frac{A^3}{S^3}+\BigO(S^{-5}).
\]
Substituting these two estimates once more into
\[
 t=\frac AS+\frac{C}{S}t^2+\frac{D}{S}t^3
 +\BigO(t^4/S)
\]
gives \eqref{eq:uq-reciprocal-reversion}.  This purely algebraic bootstrap
does not differentiate the asymptotic remainder.
\end{proof}

\begin{remark}
The theorem deliberately excludes logarithmic corrections.  If the forward
expansion contains \(t\log t\), substitution of \(t\sim A/S\) generally
produces a \((\log S)/S^3\) term in the inverse.  Such a term must not be
hidden inside an ordinary power-series ansatz.
\end{remark}

\begin{theorem}[Lambert-\(W\) preconditioner with branch compatibility]
\label{thm:uq-lambert-preconditioner}
Fix a simply connected domain \(D\subset\ComplexNumbers\setminus\{0\}\) and a holomorphic
branch of \(\log t\) on \(D\).  Consider the truncated equation
\begin{equation}
\label{eq:uq-lambert-model}
 \Lambda=\frac At+B\log t+C,
 \qquad A\ne0,\quad t\in D.
\end{equation}
If \(B=0\), its solution is \(t=A/(\Lambda-C)\), provided the denominator
is nonzero and this value lies in \(D\).  If \(B\ne0\), put
\[
 Z=-\frac AB\exp\!\left(-\frac{\Lambda-C}{B}\right).
\]
For every branch \(W_k\) defined at \(Z\) with \(W_k(Z)\ne0\), set
\[
 t_k=-\frac{A}{B\,W_k(Z)}.
\]
Then \(t_k\) solves \eqref{eq:uq-lambert-model} if and only if
\(t_k\in D\) and the branch-compatibility identity
\begin{equation}
\label{eq:uq-lambert-compatibility}
 \log t_k=\frac{\Lambda-C}{B}+W_k(Z)
\end{equation}
holds for the chosen logarithm.  Conversely, every solution of
\eqref{eq:uq-lambert-model} has this form for some value of the multivalued
relation \(u e^u=Z\), and hence for a local branch of \(W\) whenever that
relation is locally regular.
\end{theorem}

\begin{proof}
The case \(B=0\) is immediate.  Suppose \(B\ne0\) and set
\[
 u=-\frac{A}{Bt}.
\]
If \(t\) solves \eqref{eq:uq-lambert-model}, then
\[
 \log t=\frac{\Lambda-C}{B}+u.
\]
Exponentiating and using \(t=-A/(Bu)\) gives
\[
 ue^u=-\frac AB\exp\!\left(-\frac{\Lambda-C}{B}\right)=Z.
\]
Hence \(u\) is a value of \(W(Z)\); wherever \(1+u\ne0\), it belongs to a
locally holomorphic branch \(W_k\).  The displayed formula for \(t_k\) follows,
and
\eqref{eq:uq-lambert-compatibility} is the original, unexponentiated
equation.  Conversely, if \(u=W_k(Z)\), \(t_k=-A/(Bu)\), and
\eqref{eq:uq-lambert-compatibility} holds, then
\[
 \frac{A}{t_k}+B\log t_k+C
 =-Bu+B\left(\frac{\Lambda-C}{B}+u\right)+C
 =\Lambda.
\]
\end{proof}

\begin{warning}
The compatibility test in
\eqref{eq:uq-lambert-compatibility} is essential over \(\ComplexNumbers\).
Exponentiation forgets integer multiples of \(2\pi\ImaginaryUnit\); choosing a
\(W\)-branch without also fixing and checking the logarithm branch can
therefore create spurious solutions or omit genuine ones.
\end{warning}

\begin{corollary}[Accuracy of the Lambert preconditioner]
\label{cor:uq-lambert-preconditioner-error}
On a fixed real or complex branch, suppose the full equation is
\[
 \Lambda=\frac At+B\log t+C+R(t),
 \qquad A\ne0,\qquad R(t)=\BigO(t)
\]
uniformly in the chosen branch domain.  Let \(t_0\to0\) be a
branch-compatible solution of the truncated equation
\eqref{eq:uq-lambert-model}, and let \(t\) solve the full equation with
\(t/t_0\to1\).  In the complex case assume the segment joining \(t_0\) and
\(t\) stays in the chosen logarithm domain.  Then
\[
 t=t_0+\BigO(t_0^3).
\]
\end{corollary}

\begin{proof}
Write \(L(t)=A/t+B\log t+C\).  The two equations give
\[
 L(t)-L(t_0)=-R(t)=\BigO(t_0).
\]
On the joining segment, the assumption \(t/t_0\to1\) gives uniformly
\[
 L'(t)=-\frac{A}{t^2}+\frac Bt
 =-\frac{A}{t_0^2}(1+\LittleO(1)).
\]
The real mean-value theorem, or integration of \(L'\) along the complex
segment, therefore gives
\[
 L(t)-L(t_0)
 =(t-t_0)\left(-\frac{A}{t_0^2}\right)(1+\LittleO(1)).
\]
Since its left-hand side is \(\BigO(t_0)\), division proves
\(t-t_0=\BigO(t_0^3)\).
\end{proof}

\section{Multivariate triangular inversion}
\label{sec:uq-multivariate-recursion}

\begin{theorem}[Total-degree inverse recurrence]
\label{thm:uq-multivariate-triangular-inverse}
Let \(R\) be a commutative ring and let
\[
 \bm F(\bm p)=A\bm p+\sum_{r\ge2}\bm F_r(\bm p)
\]
be a \(d\)-component formal map, where each \(\bm F_r\) is homogeneous of
total degree \(r\) and \(A\in R^{d\times d}\) is invertible.  There is a
unique two-sided formal inverse
\(\bm G(\bm y)=\sum_{n\ge1}\bm G_n(\bm y)\), with each \(\bm G_n\)
homogeneous of degree \(n\).  It is determined by
\[
 \bm G_1=A^{-1}\bm y
\]
and, for \(n\ge2\),
\begin{equation}
\label{eq:uq-multivariate-recurrence}
 \bm G_n
 =-A^{-1}[\text{total degree }n]\,
 \sum_{r=2}^{n}
 \bm F_r\!\left(\sum_{j=1}^{n-1}\bm G_j\right).
\end{equation}
If \(\bm F\) is analytic over \(\RealNumbers\) or \(\ComplexNumbers\), then after translating a
base point \(\bm p_0\) and its value to the origin, the hypothesis
\(\det D\bm F(\bm p_0)\ne0\) makes the formal series the Taylor series of
the local analytic inverse, and
\[
 D\bm G(\bm y)=D\bm F(\bm G(\bm y))^{-1}.
\]
\end{theorem}

\begin{proof}
In total degree one, the equation
\(\bm F(\bm G(\bm y))=\bm y\) is \(A\bm G_1=\bm y\).
At degree \(n\ge2\), the only occurrence of the unknown
\(\bm G_n\) is the linear term \(A\bm G_n\): every nonlinear
\(\bm F_r\) uses at least two positive-degree inputs, so its degree-\(n\)
part involves only \(\bm G_1,\ldots,\bm G_{n-1}\).  Isolating
\(\bm G_n\) gives \eqref{eq:uq-multivariate-recurrence} and constructs a
unique right inverse.

For two-sidedness, let \(\bm H=\bm G\circ\bm F\).  Then
\(\bm F\circ\bm H=\bm F\) and \(D\bm H(0)=I\).  If
\(\bm H-\operatorname{id}\) first has a nonzero homogeneous part
\(\bm K_n\), then \(\bm F\circ\bm H-\bm F\) first has homogeneous part
\(A\bm K_n\), which cannot vanish because \(A\) is invertible.  Thus
\(\bm H=\operatorname{id}\).

The analytic inverse-function theorem proves convergence in the analytic
case.  Differentiating
\(\bm F(\bm G(\bm y))=\bm y\) gives
\(D\bm F(\bm G(\bm y))D\bm G(\bm y)=I\), and hence the derivative formula.
\end{proof}

\begin{remark}
When \(\det D\bm F=0\), the regular multivariate theorem is simply
inapplicable.  Rank loss alone does not identify a ramification type; a
separate normal-form or elimination argument is required.
\end{remark}

\section{The product-form Lagrange--Good formula}
\label{sec:uq-lagrange-good}

\begin{lemma}[Multivariate formal residue substitution]
\label{lem:uq-multivariate-residue}
Let \(R\) be a commutative \(\RationalNumbers\)-algebra and
\[
 T_i(\bm x)=x_iU_i(\bm x),\qquad U_i(\bm0)\in R^\times,
 \quad 1\le i\le d.
\]
For a formal Laurent series \(A(\bm T)\), whose support is bounded below in
each coordinate, define
\[
 \operatorname{Res}_{\bm T}
 A(\bm T)\,\Differential T_1\cdots\Differential T_d
 :=[T_1^{-1}\cdots T_d^{-1}]A(\bm T).
\]
Then
\[
 \operatorname{Res}_{\bm T}A(\bm T)\,\Differential T_1\cdots\Differential T_d
 =
 \operatorname{Res}_{\bm x}
 A(\bm T(\bm x))
 \det\!\left(\frac{\partial T_i}{\partial x_j}\right)
 \Differential x_1\cdots\Differential x_d.
\]
\end{lemma}

\begin{proof}
By linearity it is enough to consider a Laurent monomial
\(\prod_iT_i^{k_i}\).  If some \(k_i\ne-1\), its associated top form is
exact:
\[
 \left(\prod_jT_j^{k_j}\right)
 \Differential T_1\wedge\cdots\wedge\Differential T_d
 =
 \Differential\!\left(
 \frac{(-1)^{i-1}T_i^{k_i+1}}{k_i+1}
 \prod_{j\ne i}T_j^{k_j}
 \bigwedge_{j\ne i}\Differential T_j\right).
\]
Formal differentiation has zero residue, and pullback commutes with
\(\Differential\), so both residues vanish.  If every \(k_i=-1\), the form is
\[
 \bigwedge_{i=1}^{d}\frac{\Differential T_i}{T_i}
 =
 \bigwedge_{i=1}^{d}\left(\frac{\Differential x_i}{x_i}
 +\frac{\Differential U_i}{U_i}\right).
\]
The product of all \(\Differential x_i/x_i\) has residue one.  Every other wedge term
replaces at least one logarithmic pole by a regular power-series
differential, so its full \(x_1^{-1}\cdots x_d^{-1}\) residue is zero.
Thus the pullback also has residue one, proving the assertion.
\end{proof}

\begin{theorem}[Lagrange--Good inversion]
\label{thm:uq-lagrange-good}
Let \(R\) be a commutative \(\RationalNumbers\)-algebra and let
\(\phi_i(\bm x)\in R[[x_1,\ldots,x_d]]\) have invertible constant terms.
The product-form system
\[
 w_i=t_i\phi_i(\bm w),\qquad 1\le i\le d,
\]
has a unique formal solution \(\bm w(\bm t)\) with zero constant term.  For
every \(H(\bm x)\in R[[x_1,\ldots,x_d]]\),
\begin{equation}
\label{eq:uq-lagrange-good}
 H(\bm w(\bm t))
 =
 \sum_{\bm n\in\IntegerNumbers_{\ge0}^d}\bm t^{\bm n}
 [\bm x^{\bm n}]
 H(\bm x)\prod_{i=1}^{d}\phi_i(\bm x)^{n_i}
 \det\!\left(
 \delta_{ij}
 -\frac{x_j}{\phi_i(\bm x)}
 \frac{\partial\phi_i}{\partial x_j}(\bm x)
 \right)_{i,j=1}^{d}.
\end{equation}
\end{theorem}

\begin{proof}
Existence and uniqueness follow from
\cref{thm:uq-multivariate-triangular-inverse} applied to
\(t_i=w_i/\phi_i(\bm w)\).  Fix \(\bm n\in\IntegerNumbers_{\ge0}^d\).  Coefficient extraction
as a formal residue gives
\[
 [\bm t^{\bm n}]H(\bm w(\bm t))
 =
 \operatorname{Res}_{\bm t}
 H(\bm w(\bm t))
 \prod_i t_i^{-n_i}\frac{\Differential t_i}{t_i}.
\]
Use the formal coordinate change
\[
 t_i=\frac{x_i}{\phi_i(\bm x)}.
\]
Since \(\bm w(\bm t(\bm x))=\bm x\), one has
\[
 t_i^{-n_i}=x_i^{-n_i}\phi_i(\bm x)^{n_i}
\]
and
\[
 \frac{\Differential t_i}{t_i}
 =\sum_{j=1}^{d}\left(
 \delta_{ij}
 -\frac{x_j}{\phi_i(\bm x)}
 \frac{\partial\phi_i}{\partial x_j}(\bm x)
 \right)\frac{\Differential x_j}{x_j}.
\]
Taking the wedge product produces the determinant in
\eqref{eq:uq-lagrange-good}; taking the remaining residue is exactly
\([\bm x^{\bm n}]\).  Summing over \(\bm n\) proves the formula.
\end{proof}

\begin{remark}[Equivalent determinant conventions]
\label{rem:uq-good-determinant}
If \(D=\DiagonalOperator(x_1,\ldots,x_d)\) and
\(Q_{ij}=\phi_i^{-1}\partial\phi_i/\partial x_j\), the determinant in
\eqref{eq:uq-lagrange-good} is \(\det(I-QD)\).  Some sources instead print
\(\det(I-DQ)\), whose entries contain \(x_i\) rather than \(x_j\).
The two are equal because \(\det(I-AB)=\det(I-BA)\) for square matrices over
a commutative ring.  For completeness, this identity follows by taking the
determinant after reducing
\[
 \begin{pmatrix}I&A\\ B&I\end{pmatrix}
\]
in two different ways: the operation
\(\text{row}_2\leftarrow\text{row}_2-B\text{row}_1\) gives diagonal blocks
\((I,I-BA)\), while
\(\text{row}_1\leftarrow\text{row}_1-A\text{row}_2\) gives diagonal blocks
\((I-AB,I)\).  Both block row operations have determinant one.
\end{remark}

\begin{remark}[Scope of Good's determinant]
\label{rem:uq-good-scope}
\Cref{thm:uq-lagrange-good} applies to systems already written in the
componentwise product form \(w_i=t_i\phi_i(\bm w)\).  It is not a
coefficient formula for an arbitrary locally invertible square system.
For such a system one should first use the invariant total-degree recurrence
\eqref{eq:uq-multivariate-recurrence}; a reduction to product form requires
additional structure and must be justified.
\end{remark}

\section{Residual, Rouch\'e, and interval-Newton certificates}
\label{sec:uq-certificates}

\begin{proposition}[Real residual certificate]
\label{prop:uq-real-residual-certificate}
Let \(F\) be differentiable on an interval \(I\), suppose
\(\lvert F'(x)\rvert\ge m>0\) throughout \(I\), and let \(x_*\in I\) satisfy
\(F(x_*)=y\).  Then \(x_*\) is the unique solution of \(F(x)=y\) in \(I\),
and every \(\widetilde x\in I\) satisfies
\[
 \lvert\widetilde x-x_*\rvert
 \le\frac{\lvert F(\widetilde x)-y\rvert}{m}.
\]
\end{proposition}

\begin{proof}
Derivatives have the Darboux property, so a derivative bounded away from zero
cannot change sign; hence \(F\) is strictly monotone on \(I\) and the root is
unique.  Apply \cref{prop:uq-asymptotic-residual-transfer} with
\(G(y)=x_*\), \(\widetilde G(y)=\widetilde x\), and \(m(y)=m\); its segment
hypothesis follows from the bound on all of \(I\).
\end{proof}

\begin{theorem}[Rouch\'e disk certificate]
\label{thm:uq-rouche-disk-certificate}
Let \(r>0\), let \(F\) be holomorphic on a neighbourhood of the closed disk
\(\overline D(z_0,r)\), let \(y\in\ComplexNumbers\), and assume
\[
 F'(z_0)\ne0,\qquad
 \sup_{\overline D(z_0,r)}\lvert F''(z)\rvert\le M.
\]
If
\begin{equation}
\label{eq:uq-rouche-condition}
 \frac{Mr^2}{2}
 <
 \lvert F'(z_0)\rvert r-\lvert F(z_0)-y\rvert,
\end{equation}
then \(F(z)=y\) has exactly one solution in \(D(z_0,r)\), and that solution
is simple.
\end{theorem}

\begin{proof}
Put \(H(z)=F(z)-y\) and
\[
 L(z)=H(z_0)+F'(z_0)(z-z_0).
\]
Taylor's formula along the segment from \(z_0\) to \(z\), valid because a
disk is convex, gives
\[
 \lvert H(z)-L(z)\rvert\le\frac{M}{2}\lvert z-z_0\rvert^2.
\]
On the boundary \(\lvert z-z_0\rvert=r\),
\[
 \lvert L(z)\rvert
 \ge\lvert F'(z_0)\rvert r-\lvert H(z_0)\rvert
 >\frac{Mr^2}{2}
 \ge\lvert H(z)-L(z)\rvert.
\]
Rouch\'e's theorem says that \(H\) and \(L\) have the same number of zeros
in the disk, counted with multiplicity.  Condition
\eqref{eq:uq-rouche-condition} also implies
\(\lvert H(z_0)/F'(z_0)\rvert<r\), so the affine function \(L\) has exactly
one zero there.  Hence \(H\) has total zero multiplicity one, which means
its unique zero is simple.
\end{proof}

\begin{theorem}[Scalar interval-Newton certificate]
\label{thm:uq-interval-newton-certificate}
Let \(I=[a,b]\subset\RealNumbers\), let \(F\in C^1(I)\), choose \(x_0\in I\), and let
\(D=[d_-,d_+]\) be a closed real interval containing \(F'(I)\), with
\(0\notin D\).  Interpret \(F(x_0)/D\) as its interval image and define
\[
 N=x_0-\frac{F(x_0)}{D}.
\]
If \(N\subseteq I\), then \(F\) has exactly one zero \(x_*\) in \(I\), and
\(x_*\in N\).  If \(N\subseteq\operatorname{int}I\), the certified zero lies
in \(\operatorname{int}I\).
\end{theorem}

\begin{proof}
Replacing \(F\) by \(-F\) and \(D\) by \(-D\), if necessary, leaves \(N\)
unchanged, so assume
\(0<d_-\le F'(x)\le d_+\) on \(I\).  Thus \(F\) is strictly increasing.
If \(F(x_0)=0\), the conclusion is immediate.  If \(F(x_0)>0\), then
\[
 N=\left[
 x_0-\frac{F(x_0)}{d_-},
 x_0-\frac{F(x_0)}{d_+}
 \right]\subseteq[a,b].
\]
The left-end inclusion implies
\(F(x_0)\le d_-(x_0-a)\).  Since
\[
 F(x_0)-F(a)=\int_a^{x_0}F'(s)\,\Differential s
 \ge d_-(x_0-a),
\]
we obtain \(F(a)\le0<F(x_0)\), so the intermediate-value theorem gives a
zero between \(a\) and \(x_0\).  If \(F(x_0)<0\), the right-end inclusion
similarly gives
\(-F(x_0)\le d_-(b-x_0)\), hence \(F(b)\ge0\), and again a zero exists.
Strict monotonicity makes it unique.

For that zero \(x_*\), the mean-value theorem gives some \(\theta\) between
\(x_0\) and \(x_*\) such that
\[
 F(x_0)=F'(\theta)(x_0-x_*).
\]
Because \(F'(\theta)\in D\),
\[
 x_*=x_0-\frac{F(x_0)}{F'(\theta)}\in N.
\]
The final interior assertion follows from \(x_*\in N\).
\end{proof}

\begin{boxedremark}
These certificates answer different logical questions.  A residual estimate
controls the error once existence and the intended branch are already known.
The Rouch\'e test proves complex existence, uniqueness, and simplicity in a
disk.  The scalar interval-Newton test proves real existence, uniqueness, and
localization in an interval.  A small numerical residual alone must never be
reported as an existence or branch certificate.
\end{boxedremark}
